%% Cover letter for the submission of "A Littlewood--Richardson rule for forest
%% polynomials" to Annals of Combinatorics.  Build: pdflatex cover_letter
\documentclass[11pt]{article}
\usepackage[margin=1in]{geometry}
\usepackage{amsmath,amssymb}
\usepackage[hidelinks]{hyperref}
\pagestyle{empty}
\setlength{\parindent}{0pt}
\setlength{\parskip}{0.6\baselineskip}

\begin{document}

Matthew J. Samuel\\
Monroe Township, Middlesex County, NJ, USA\\
\href{mailto:matthematics@gmail.com}{matthematics@gmail.com} \quad$\cdot$\quad ORCID \href{https://orcid.org/0009-0000-1833-0890}{0009-0000-1833-0890}

October 28, 2026

The Editors\\
\emph{Annals of Combinatorics}

Dear Editors,

I am pleased to submit the manuscript \emph{A Littlewood--Richardson rule for forest polynomials} for consideration in \emph{Annals of Combinatorics}, under the Algebraic Combinatorics track.

The forest polynomials of Nadeau and Tewari form a basis of the polynomial ring that plays, for the quasisymmetric flag variety of Bergeron, Gagnon, Nadeau, Spink and Tewari, the role played by Schubert polynomials for the classical flag variety. Their structure constants are known to be nonnegative, but no Littlewood--Richardson-type rule for them has been available. This paper gives such a rule: each coefficient is the number of pairs of RC graphs (pipe dreams) with prescribed codes whose ``lift product'' is a distinguished RC graph of prescribed weight, and the same count computes cup products on the quasisymmetric flag variety. The proof introduces products on RC graphs that recover the plactic monoid, model Schubert polynomial multiplication, and carry a quasi-crystal structure, which I believe are of independent interest. Combining Schubert calculus, symmetric functions, crystals, and the combinatorics of reduced words, the paper falls squarely within the journal's Algebraic Combinatorics track. Some years ago Nathan Williams kindly encouraged me to submit work to the journal, and I am glad to be able to do so now; given the subject matter, Vasu Tewari may be a natural choice of handling editor.

The manuscript is a companion to \emph{The Schubert bialgebra} (arXiv:2610.22582), which has been submitted to another journal; the present paper is self-contained, collecting with precise references everything it uses from the companion. For transparency, the first arXiv version of this work (arXiv:2606.23876v1) contained the material of both papers, which has since been split. The manuscript is my own original work, has not been published, and is not under consideration elsewhere; I am its sole author, have no competing interests, and received no funding for this work. I am an independent researcher, not affiliated with an institution.

Thank you for your consideration.

Sincerely,

\vspace{1.2\baselineskip}
Matthew J. Samuel

\end{document}
